\documentclass[conference]{IEEEtran}
\IEEEoverridecommandlockouts

\usepackage{cite}
\usepackage{amsmath,amssymb,amsfonts}
\usepackage{graphicx}
\usepackage{textcomp}
\usepackage{tabularx,booktabs}
\usepackage{xcolor}
\usepackage{hyperref}
\usepackage{url}

\begin{document}

\title{Facilitando a Criação de Experiências Educacionais em Realidade Aumentada por Usuários Não Técnicos:  
Um Mapeamento Exploratório e Avaliação Técnica de Ferramentas de IA e WebAR}

\author{
\IEEEauthorblockN{
Jorge Luiz Cunha de Freitas
}
\IEEEauthorblockA{
Centro de Informática (CIn)\\
Universidade Federal de Pernambuco (UFPE)\\
Recife, Brasil\\
Email: jlf@cin.ufpe.br
}
}

\maketitle

% ---------------------------
% RESUMO
% ---------------------------
\begin{abstract}
A Realidade Aumentada (RA) tem se consolidado como uma tecnologia promissora para aplicações educacionais, especialmente por seu potencial de promover experiências interativas e multimodais. Entretanto, a criação de experiências personalizadas em RA ainda impõe barreiras técnicas significativas a usuários sem experiência prévia em modelagem 3D, desenvolvimento de software ou uso de motores gráficos complexos. Paralelamente, ferramentas baseadas em Inteligência Artificial (IA) têm emergido com o objetivo de democratizar a criação de modelos tridimensionais, permitindo a geração de avatares e objetos 3D a partir de interfaces simplificadas e descrições textuais.

Este trabalho investiga a viabilidade de usuários não técnicos criarem experiências educacionais completas em RA, combinando ferramentas de IA para geração de modelos 3D com tecnologias WebAR acessíveis diretamente via navegador. A pesquisa adota um mapeamento exploratório das principais plataformas de criação de modelos 3D baseadas em IA, seguido de um benchmark técnico comparativo, no qual são avaliados critérios como facilidade de uso, qualidade visual, otimização das malhas, tamanho dos arquivos e compatibilidade com bibliotecas WebAR.

Como contribuição prática, foi desenvolvido um fluxo experimental completo e uma plataforma WebAR de teste para validação dos modelos GLB gerados pelos usuários em dispositivos móveis. Os resultados indicam que já é possível criar experiências de RA funcionais, acessíveis e publicáveis na Web sem o uso de motores gráficos tradicionais, embora ainda exista uma lacuna significativa quanto à oferta de soluções integradas no-code voltadas a usuários iniciantes.
\end{abstract}

% ---------------------------
% ABSTRACT
% ---------------------------
\begin{abstract}
Augmented Reality (AR) has emerged as a promising technology for educational applications, particularly due to its potential to promote interactive and multimodal learning experiences. However, the creation of personalized AR content still poses significant technical barriers for users without prior experience in 3D modeling, software development, or complex game engines. In parallel, Artificial Intelligence (AI)-based tools have emerged with the goal of democratizing 3D content creation, enabling the generation of avatars and 3D objects through simplified interfaces and textual descriptions.

This work investigates the feasibility of non-technical users creating complete educational AR experiences by combining AI-based 3D model generation tools with WebAR technologies accessible directly through web browsers. The study adopts an exploratory mapping of major AI-driven 3D creation platforms, followed by a comparative technical benchmark evaluating criteria such as usability, visual quality, mesh optimization, file size, and compatibility with WebAR libraries.

As a practical contribution, a complete experimental workflow and a WebAR testing platform were developed to validate GLB models generated by users on mobile devices. The results indicate that it is already possible to create functional, accessible, and web-publishable AR experiences without relying on traditional game engines, although a significant gap remains in the availability of integrated no-code solutions tailored to beginner users.
\end{abstract}


\begin{IEEEkeywords}
Realidade Aumentada, WebAR, Inteligência Artificial, Criação 3D, GLB, Acessibilidade Tecnológica
\end{IEEEkeywords}

% ---------------------------
% INTRODUÇÃO
% ---------------------------
\section{Introdução}

A Realidade Aumentada (RA) evoluiu significativamente desde as definições clássicas apresentadas por Azuma \cite{azuma1997} e pela taxonomia de Realidade Misturada proposta por Milgram e Kishino \cite{milgram1994}. Originalmente associada a aplicações industriais e científicas, a RA passou a ser amplamente explorada em contextos educacionais, culturais e sociais, em função do aumento da capacidade computacional de dispositivos móveis e da popularização de tecnologias gráficas acessíveis.

Avanços recentes em tecnologias baseadas na Web, como WebXR \cite{webxr2020}, possibilitaram a execução de experiências imersivas diretamente no navegador, ampliando o acesso a aplicações de RA sem a necessidade de instalação de aplicativos nativos. Esse paradigma, conhecido como WebAR, tem sido amplamente estudado por reduzir barreiras técnicas e facilitar a disseminação de conteúdos interativos em larga escala.

Nesse contexto, bibliotecas WebAR como AR.js \cite{arjs} tornaram-se populares por permitirem experiências baseadas em marcadores com baixo custo computacional, sendo particularmente adequadas para dispositivos móveis com recursos limitados. Tais características tornam o WebAR especialmente atrativo para aplicações educacionais, nas quais simplicidade de acesso e compatibilidade multiplataforma são requisitos fundamentais.

Paralelamente, ferramentas de Inteligência Artificial (IA) voltadas à geração automática de modelos 3D evoluíram rapidamente. Plataformas como Meshy AI \cite{meshy2024}, Luma AI \cite{lumaai2024} e Ready Player Me \cite{readyplayerme} permitem que usuários criem avatares e objetos tridimensionais a partir de descrições textuais ou interfaces guiadas, reduzindo significativamente a necessidade de conhecimentos avançados em modelagem 3D. Estudos recentes indicam que essas ferramentas contribuem para a democratização da criação de conteúdo tridimensional, ampliando o acesso a tecnologias antes restritas a especialistas.

Apesar desses avanços, integrar modelos 3D personalizados em experiências WebAR ainda representa um desafio considerável para usuários não técnicos. Problemas relacionados à compatibilidade de formatos, tamanho excessivo dos arquivos, ausência de otimização de malhas e a inexistência de plataformas integradas que unifiquem geração, validação e publicação dificultam a criação de fluxos completos de autoria \cite{ruiz2020arbeginner}. Como consequência, muitos usuários dependem de múltiplas ferramentas independentes ou abandonam o processo antes de concluir uma experiência funcional.

Diante desse cenário, este trabalho investiga a seguinte questão central: \textit{até que ponto usuários sem conhecimento técnico conseguem criar experiências educacionais completas em RA utilizando ferramentas modernas de IA e tecnologias WebAR?} Para responder a essa questão, o estudo combina uma revisão exploratória das principais plataformas disponíveis com um benchmark técnico comparativo e a implementação de uma plataforma WebAR experimental para validação prática dos modelos gerados.

% ---------------------------
% TRABALHOS RELACIONADOS
% ---------------------------
\section{Trabalhos Relacionados}

A Realidade Aumentada baseada na Web (WebAR) tem sido amplamente investigada como alternativa viável para aplicações educacionais, sobretudo por eliminar a necessidade de instalação de aplicativos nativos e reduzir barreiras de acesso tecnológico. Estudos como o de Koutsabasis et al. \cite{koutsabasis2021webar} destacam que soluções WebAR favorecem a adoção em contextos educacionais formais e informais, especialmente quando combinadas com dispositivos móveis amplamente disponíveis.

Revisões sistemáticas da literatura apontam que o uso de RA no ensino pode resultar em ganhos significativos de engajamento, motivação e compreensão conceitual, particularmente em disciplinas que se beneficiam de visualização espacial e interação tridimensional \cite{fonnet2021review, munoz2021education}. No entanto, esses estudos também ressaltam que grande parte das aplicações analisadas são desenvolvidas por equipes técnicas especializadas, limitando a replicabilidade e a personalização por professores e alunos.

No que diz respeito à autoria de conteúdo, pesquisas indicam que usuários iniciantes enfrentam dificuldades substanciais na criação, preparação e publicação de modelos 3D para RA. Ruiz et al. \cite{ruiz2020arbeginner} demonstram que desafios como otimização de malhas, controle de tamanho dos arquivos e compatibilidade entre formatos representam barreiras recorrentes, especialmente em ambientes WebAR voltados a dispositivos móveis com recursos limitados.

Paralelamente, trabalhos recentes sobre geração tridimensional assistida por Inteligência Artificial evidenciam avanços significativos na qualidade visual e na automação do processo de modelagem. Plataformas baseadas em IA têm sido estudadas como mecanismos de democratização da criação 3D, permitindo a geração de objetos e avatares a partir de descrições textuais ou imagens de referência \cite{mesh2023}. Apesar disso, esses estudos frequentemente abordam a geração de modelos de forma isolada, sem considerar sua integração em fluxos completos de publicação e execução em RA.

Além disso, a literatura aponta uma lacuna quanto à avaliação sistemática da compatibilidade entre modelos gerados por IA e tecnologias WebAR. Poucos trabalhos realizam análises técnicas envolvendo métricas como tamanho de arquivos, contagem de polígonos e desempenho em navegadores móveis, fatores críticos para experiências WebAR funcionais.

Diferentemente dos trabalhos existentes, esta pesquisa adota uma abordagem prática e comparativa, avaliando plataformas de geração 3D com foco explícito na compatibilidade com WebAR e na perspectiva de usuários não técnicos. Ademais, o estudo propõe e implementa um fluxo experimental completo, incluindo uma plataforma WebAR de validação para testes diretos dos modelos GLB gerados, contribuindo para a compreensão dos limites e das possibilidades atuais de criação de experiências educacionais em RA de forma acessível.

% ---------------------------
% METODOLOGIA
% ---------------------------
\section{Metodologia}

A metodologia adotada combina um mapeamento exploratório de ferramentas com um benchmark técnico comparativo e experimentação prática. A abordagem foi estruturada com o objetivo de avaliar, sob a perspectiva de usuários não técnicos, a viabilidade de criação de experiências educacionais em Realidade Aumentada utilizando ferramentas de Inteligência Artificial e tecnologias WebAR. O processo metodológico foi dividido em três etapas principais:

\begin{enumerate}
    \item levantamento e seleção de plataformas de criação 3D baseadas em IA;
    \item definição de critérios técnicos de avaliação alinhados às restrições de WebAR;
    \item implementação de um fluxo experimental e de uma plataforma WebAR para validação prática dos modelos gerados.
\end{enumerate}

\subsection{Mapeamento Exploratório de Ferramentas}

O mapeamento exploratório teve como objetivo identificar plataformas de geração de modelos 3D baseadas em IA com potencial de uso por usuários iniciantes. Para isso, foram realizadas buscas em catálogos especializados de ferramentas de IA, como agregadores públicos de soluções emergentes, além da análise de documentação técnica oficial e artigos científicos recentes.

Os critérios iniciais de triagem priorizaram ferramentas que:
\begin{itemize}
    \item fossem acessíveis via navegador, sem necessidade de instalação local;
    \item oferecessem fluxo de criação simplificado, adequado a usuários sem experiência em modelagem 3D;
    \item permitissem exportação direta no formato GLB;
    \item apresentassem evidências de uso em contextos educacionais ou de prototipação rápida.
\end{itemize}

Plataformas que exigiam conhecimento avançado de modelagem, ferramentas profissionais complexas ou pipelines de exportação pouco documentados foram descartadas nesta fase.

\subsection{Plataformas Avaliadas}

A partir do mapeamento exploratório, foram selecionadas três plataformas representativas, consideradas relevantes tanto pela popularidade quanto pela diversidade de abordagens de geração 3D:

\begin{itemize}
    \item \textbf{Ready Player Me}: plataforma voltada à criação de avatares humanos estilizados, com forte foco em acessibilidade e uso em aplicações interativas.
    \item \textbf{Meshy AI}: ferramenta baseada em IA generativa para criação de modelos 3D a partir de descrições textuais ou imagens, com ênfase em otimização automática.
    \item \textbf{Luma AI (Genie)}: solução baseada em geração neural de objetos tridimensionais, priorizando fidelidade visual e realismo.
\end{itemize}

A seleção dessas plataformas permitiu comparar diferentes paradigmas de geração 3D, incluindo avatares humanos e objetos genéricos, bem como diferentes níveis de controle e automação.

\subsection{Definição dos Critérios de Avaliação}

Os critérios de avaliação foram definidos com base em recomendações da literatura sobre WebAR e em limitações práticas observadas em dispositivos móveis \cite{koutsabasis2021webar, ruiz2020arbeginner}. Sete critérios técnicos foram adotados:

\begin{itemize}
    \item \textbf{Facilidade de uso}: número de etapas, clareza da interface e necessidade de decisões técnicas por parte do usuário;
    \item \textbf{Qualidade visual do modelo}: fidelidade estética, coerência geométrica e aparência geral;
    \item \textbf{Exportação GLB}: simplicidade e confiabilidade do processo de exportação;
    \item \textbf{Tamanho do arquivo}: peso final do modelo exportado, relevante para carregamento em WebAR;
    \item \textbf{Contagem de polígonos}: número aproximado de triângulos, impactando desempenho em navegadores móveis;
    \item \textbf{Compatibilidade com WebAR}: capacidade de execução correta em bibliotecas WebAR sem ajustes adicionais;
    \item \textbf{Tempo total do fluxo}: tempo necessário desde a criação inicial até a obtenção de um modelo funcional.
\end{itemize}

Cada critério foi pontuado em uma escala ordinal de 0 a 3, onde valores mais altos indicam melhor adequação ao contexto de WebAR, totalizando até 21 pontos por plataforma.

\subsection{Procedimento Experimental}

Para garantir comparabilidade entre as plataformas, foi definido um procedimento experimental padronizado. Em cada ferramenta, buscou-se gerar um avatar infantil com características semelhantes, utilizando descrições equivalentes sempre que aplicável.

O procedimento experimental seguiu as seguintes etapas:

\begin{enumerate}
    \item criação do modelo 3D na plataforma selecionada;
    \item exportação do modelo no formato GLB;
    \item validação estrutural do arquivo utilizando o \textit{Khronos glTF Validator};
    \item análise do tamanho do arquivo e contagem de polígonos;
    \item testes de execução em ambiente WebAR utilizando AR.js e ModelViewer;
    \item registro do tempo total do fluxo de criação.
\end{enumerate}

Todos os experimentos foram conduzidos utilizando o mesmo ambiente computacional e testados em dispositivos móveis padrão, a fim de reduzir variações externas.

\subsection{Plataforma WebAR de Validação}

Como parte da metodologia, foi desenvolvida uma plataforma WebAR leve com o objetivo de permitir a validação prática dos modelos GLB gerados. Essa plataforma possibilita o carregamento direto dos modelos em uma cena WebAR baseada em marcadores, facilitando a inspeção de desempenho, escala, orientação e estabilidade do modelo no navegador.

A plataforma foi implementada utilizando HTML, JavaScript e a biblioteca AR.js, sendo hospedada em ambiente público para acesso remoto. Essa etapa permitiu simular o uso real por usuários finais e identificar limitações que não seriam perceptíveis apenas por análise estática dos arquivos.

A inclusão dessa plataforma de validação contribui para aproximar o benchmark técnico de um cenário real de uso educacional, reforçando o caráter aplicado da pesquisa.

% ---------------------------
% PLATAFORMA WEBAR DE TESTE
% ---------------------------
\section{Plataforma WebAR para Validação dos Modelos}

Como contribuição prática deste trabalho, foi desenvolvida uma plataforma WebAR baseada na biblioteca AR.js, com o objetivo de validar, de forma padronizada e reprodutível, os modelos GLB gerados pelas ferramentas de Inteligência Artificial avaliadas. A plataforma foi concebida para simular um cenário real de uso educacional, permitindo que usuários carreguem modelos 3D personalizados diretamente no navegador, sem a necessidade de instalação de aplicativos nativos.

A escolha do AR.js deve-se à sua ampla adoção, baixo custo computacional e compatibilidade com dispositivos móveis, características frequentemente destacadas na literatura como fundamentais para aplicações WebAR acessíveis \cite{koutsabasis2021webar}. A plataforma utiliza RA baseada em marcadores, abordagem que apresenta maior estabilidade e previsibilidade para usuários iniciantes, além de facilitar a criação de materiais educacionais impressos, como cartilhas e pôsteres.

A arquitetura da plataforma foi implementada utilizando tecnologias web padrão (HTML, JavaScript e A-Frame), permitindo o carregamento dinâmico de modelos no formato GLB. Cada modelo é posicionado automaticamente sobre um marcador personalizado, garantindo consistência na escala, orientação e ponto de ancoragem durante os testes. Essa padronização foi essencial para comparar o comportamento dos modelos gerados por diferentes plataformas em condições equivalentes.

Além da visualização tridimensional, a plataforma suporta a reprodução de áudio sintetizado por IA, associado à presença do marcador. Esse recurso possibilita a criação de experiências multimodais, combinando elementos visuais e sonoros, frequentemente apontadas como eficazes em contextos educacionais baseados em RA \cite{fonnet2021review}. O áudio é acionado automaticamente após o reconhecimento do marcador, simulando narrativas educativas conduzidas por personagens virtuais.

A plataforma foi hospedada em ambiente público com suporte a HTTPS, permitindo acesso remoto e testes em diferentes dispositivos móveis. Essa decisão viabilizou a avaliação do desempenho em navegadores reais, considerando fatores como tempo de carregamento, estabilidade da renderização e responsividade da interação.

O desenvolvimento dessa plataforma mostrou-se fundamental para avaliar não apenas a qualidade visual dos modelos gerados, mas principalmente sua viabilidade prática de uso em cenários educacionais reais. Problemas como arquivos excessivamente pesados, malhas não otimizadas ou incompatibilidades de materiais tornaram-se evidentes apenas durante a execução em WebAR, reforçando a importância de uma validação prática integrada ao benchmark técnico.

Por fim, a plataforma WebAR desenvolvida também aponta para uma linha futura de pesquisa: a criação de uma solução integrada no-code, na qual usuários finais possam gerar seus próprios avatares por IA, associá-los a marcadores personalizados e publicar experiências educacionais completas sem conhecimento técnico especializado.


% ---------------------------
% RESULTADOS
% ---------------------------
\section{Resultados}

Os resultados obtidos a partir do benchmark técnico e da validação prática em WebAR indicam que todas as plataformas avaliadas permitem a criação de modelos 3D compatíveis com ambientes WebAR. Entretanto, foram observadas diferenças significativas em termos de desempenho técnico, otimização para dispositivos móveis e adequação ao uso por usuários não técnicos.

\subsection{Análise Qualitativa dos Resultados}

Além dos indicadores quantitativos obtidos no benchmark técnico, foi realizada uma análise qualitativa do processo de criação e integração dos modelos 3D, considerando a perspectiva de um usuário não técnico. Observou-se que a clareza das interfaces e o nível de abstração oferecido pelas plataformas influenciam diretamente a percepção de dificuldade do fluxo de criação, independentemente da qualidade final do modelo gerado.

No caso do Ready Player Me, a experiência de criação mostrou-se altamente guiada e intuitiva, com etapas bem definidas e feedback visual constante, o que reduz significativamente a carga cognitiva do usuário iniciante. Já o Meshy AI, embora tecnicamente superior em termos de otimização dos modelos, exige maior compreensão conceitual sobre prompts e parâmetros, o que pode representar uma barreira inicial, apesar dos excelentes resultados obtidos.

A Luma AI destacou-se pela expressividade visual dos modelos gerados, porém a ausência de controles claros sobre otimização e complexidade da malha pode levar usuários iniciantes a produzir modelos inadequados para WebAR sem perceber essas limitações durante o processo de criação. Essa diferença entre percepção visual e viabilidade técnica foi recorrente nos testes realizados.

De forma geral, a análise qualitativa evidencia que a acessibilidade percebida pelo usuário não depende exclusivamente da simplicidade da interface, mas também da capacidade da plataforma de guiar o processo de criação em direção a resultados tecnicamente compatíveis com WebAR. Esse aspecto reforça a necessidade de soluções integradas que incorporem validações automáticas, sugestões de otimização e alertas durante o processo de autoria.


\subsection{Resultados Quantitativos}

A Tabela~\ref{tab:benchmark} apresenta a pontuação consolidada das plataformas avaliadas, considerando os sete critérios definidos na metodologia. Observa-se que o Meshy AI obteve a maior pontuação geral, refletindo um equilíbrio consistente entre qualidade visual, otimização da malha e compatibilidade com WebAR.

\begin{table}[ht]
\caption{Benchmark técnico das plataformas avaliadas.}
\label{tab:benchmark}
\centering
\begin{tabular}{lcccccccc}
\toprule
\textbf{Plataforma} & A & B & C & D & E & F & G & \textbf{Total} \\
\midrule
Ready Player Me & 3 & 2 & 3 & 2 & 2 & 3 & 3 & \textbf{18} \\
Meshy AI        & 3 & 3 & 3 & 3 & 3 & 3 & 2 & \textbf{20} \\
Luma AI         & 2 & 3 & 2 & 1 & 1 & 2 & 2 & \textbf{13} \\
\bottomrule
\end{tabular}
\end{table}


O Ready Player Me apresentou desempenho elevado nos critérios de facilidade de uso, exportação GLB e tempo total do fluxo, evidenciando sua adequação para usuários iniciantes que desejam criar rapidamente personagens humanoides. Em contrapartida, seus modelos apresentaram tamanhos de arquivo e contagem de polígonos ligeiramente superiores aos do Meshy AI.

A Luma AI (Genie) destacou-se pela qualidade visual dos modelos gerados, alcançando pontuação máxima nesse critério. Contudo, apresentou os piores resultados em tamanho de arquivo e contagem de polígonos, impactando diretamente o tempo de carregamento e a estabilidade da renderização em dispositivos móveis.

\subsection{Resultados Qualitativos em Ambiente WebAR}

A validação prática realizada na plataforma WebAR desenvolvida evidenciou diferenças relevantes no comportamento dos modelos durante a execução em navegadores móveis. Modelos gerados pelo Meshy AI foram carregados de forma quase instantânea, apresentando estabilidade na detecção do marcador e fluidez na renderização, mesmo em dispositivos com menor capacidade de processamento.

Os modelos do Ready Player Me apresentaram desempenho satisfatório na maioria dos testes, embora tenham sido observados pequenos atrasos no carregamento inicial em conexões mais lentas. Ainda assim, não foram identificados erros críticos de renderização ou incompatibilidades com AR.js.

Já os modelos gerados pela Luma AI apresentaram maior latência no carregamento e, em alguns casos, warnings no validador glTF relacionados a complexidade excessiva da malha. Em dispositivos móveis mais antigos, foi observado aquecimento e redução na taxa de quadros, o que pode comprometer a experiência do usuário final em contextos educacionais.

\subsection{Análise do Fluxo de Criação por Usuários Não Técnicos}

Do ponto de vista do usuário não técnico, o fluxo de criação variou significativamente entre as plataformas. O Ready Player Me mostrou-se o mais intuitivo, permitindo a criação de um avatar funcional em poucos minutos, com mínima necessidade de ajustes adicionais. Esse fator é especialmente relevante para professores e educadores sem formação técnica.

O Meshy AI exigiu maior atenção na formulação dos prompts textuais, porém compensou com modelos mais otimizados e prontos para uso em WebAR, reduzindo a necessidade de retrabalho posterior. A Luma AI, embora simples de operar, gerou modelos que demandariam etapas adicionais de otimização, o que pode representar uma barreira para usuários iniciantes.

De forma geral, os resultados indicam que a acessibilidade percebida pelo usuário não depende apenas da facilidade de geração do modelo, mas também da compatibilidade imediata com o ambiente WebAR, reforçando a importância de considerar critérios técnicos desde as etapas iniciais do fluxo de criação.


% ---------------------------
% DISCUSSÃO
% ---------------------------
\section{Discussão}

Os resultados obtidos neste estudo confirmam que usuários não técnicos já dispõem de ferramentas suficientemente acessíveis para criar experiências de Realidade Aumentada funcionais, utilizando a combinação entre geração de modelos 3D por Inteligência Artificial e tecnologias WebAR. Esse achado está alinhado com trabalhos recentes que apontam a WebAR como alternativa viável para ampliar o acesso a aplicações de RA, especialmente em contextos educacionais \cite{koutsabasis2021webar, fonnet2021review}.

Entretanto, os experimentos também evidenciam que a acessibilidade percebida pelo usuário final não depende apenas da facilidade de geração do modelo 3D. Fatores técnicos como tamanho do arquivo, contagem de polígonos e compatibilidade imediata com bibliotecas WebAR têm impacto direto na experiência, corroborando observações feitas por Ruiz et al. \cite{ruiz2020arbeginner} sobre as dificuldades enfrentadas por iniciantes na preparação de conteúdo para RA.

A análise comparativa mostrou que plataformas como o Meshy AI oferecem modelos mais bem otimizados para execução em dispositivos móveis, enquanto soluções como Ready Player Me priorizam a simplicidade do fluxo de criação. Já ferramentas mais recentes, como a Luma AI, demonstram elevado potencial visual, porém ainda carecem de otimizações automáticas para uso em ambientes WebAR. Esses resultados reforçam a necessidade de critérios técnicos claros para avaliação de ferramentas de criação 3D voltadas a RA baseada em navegador.

A plataforma WebAR desenvolvida neste trabalho desempenhou papel central na validação prática dos modelos gerados. Ao fornecer um ambiente padronizado para carregamento, visualização e teste dos arquivos GLB, foi possível avaliar não apenas a qualidade estética dos modelos, mas sua viabilidade real de uso em cenários educacionais. Essa abordagem experimental diferencia este estudo de revisões puramente teóricas, ao evidenciar problemas que só emergem durante a execução em dispositivos móveis reais.

Além disso, a plataforma de teste demonstra, na prática, a viabilidade técnica de integrar múltiplas etapas do fluxo de autoria — geração do modelo, validação, inserção de áudio e publicação WebAR — em uma única solução baseada na Web. Embora o protótipo desenvolvido não configure ainda uma plataforma completa no-code, ele evidencia que tal integração é tecnicamente possível com tecnologias atuais, reforçando a relevância de pesquisas futuras nessa direção.

Por fim, os resultados indicam que a principal lacuna atual não está na ausência de ferramentas individuais, mas na falta de um ecossistema integrado que permita a usuários não técnicos criar, testar e publicar experiências de RA de forma contínua e sem retrabalho. Essa lacuna representa uma oportunidade significativa para o desenvolvimento de plataformas educacionais acessíveis, alinhadas às tendências recentes de democratização da criação de conteúdo digital.



% ---------------------------
% LIMITAÇÕES DO ESTUDO
% ---------------------------

\section{Limitações do Estudo}

Apesar dos resultados promissores apresentados, este estudo possui algumas limitações que devem ser consideradas na interpretação dos achados.

Primeiramente, o número de plataformas avaliadas foi restrito a três ferramentas de geração de modelos 3D baseadas em IA. Embora essas plataformas sejam representativas de abordagens distintas e amplamente utilizadas, o ecossistema atual é dinâmico e inclui outras soluções emergentes que não foram contempladas neste trabalho. Assim, os resultados não devem ser generalizados para todas as ferramentas disponíveis no mercado.

Outra limitação refere-se à variabilidade inerente aos modelos gerados por Inteligência Artificial. Pequenas alterações nos parâmetros de entrada, descrições textuais ou atualizações internas das plataformas podem resultar em diferenças significativas na qualidade, no tamanho e na complexidade dos modelos 3D produzidos. Para mitigar esse fator, buscou-se utilizar descrições e configurações semelhantes em todas as plataformas, porém a completa reprodutibilidade não pode ser garantida.

Além disso, a avaliação experimental foi conduzida em um conjunto específico de dispositivos (um notebook e um smartphone), o que pode não refletir integralmente o desempenho em outros contextos de hardware, especialmente em dispositivos móveis de menor capacidade computacional. Questões como desempenho gráfico, tempo de carregamento e estabilidade da experiência WebAR podem variar de acordo com o dispositivo utilizado.

Outro aspecto limitante está relacionado ao escopo da plataforma WebAR desenvolvida. Embora ela tenha sido projetada para validar a compatibilidade e a viabilidade prática dos modelos GLB em cenários educacionais, não foram realizadas avaliações formais com usuários finais, como professores ou estudantes. Dessa forma, aspectos de usabilidade, aceitação e impacto pedagógico não foram investigados neste estudo.

Por fim, o foco deste trabalho esteve predominantemente na viabilidade técnica e operacional da criação e integração de modelos 3D em WebAR. Elementos mais avançados de interação, como reconhecimento de gestos, múltiplos marcadores simultâneos ou sincronização labial avançada, não foram considerados, representando oportunidades para investigações futuras.



% ---------------------------
% CONCLUSÃO
% ---------------------------
\section{Conclusão}

Este trabalho apresentou um estudo exploratório sobre a criação de experiências educacionais em Realidade Aumentada por usuários não técnicos, combinando ferramentas de Inteligência Artificial para geração de modelos 3D com tecnologias WebAR acessíveis via navegador. A pesquisa abordou tanto aspectos conceituais, por meio de um mapeamento exploratório da literatura e das plataformas disponíveis, quanto aspectos práticos, através de um benchmark técnico comparativo e da implementação de uma plataforma WebAR para validação de modelos GLB.

Os resultados obtidos demonstram que já é tecnicamente viável para usuários iniciantes criarem modelos 3D compatíveis com WebAR utilizando ferramentas baseadas em IA, sem a necessidade de conhecimentos avançados em modelagem 3D ou motores gráficos. Plataformas como Meshy AI e Ready Player Me mostraram-se particularmente adequadas para esse contexto, destacando-se, respectivamente, pela otimização técnica dos modelos e pela facilidade de uso no fluxo de criação.

Entretanto, a avaliação também evidenciou limitações importantes. Diferenças significativas no tamanho dos arquivos, na contagem de polígonos e na compatibilidade imediata com bibliotecas WebAR impactam diretamente a experiência final, especialmente em dispositivos móveis. Esses resultados confirmam que a acessibilidade percebida pelo usuário final não depende apenas da interface das ferramentas de criação, mas também da adequação técnica dos modelos gerados aos requisitos da WebAR.

Como principal contribuição prática, este trabalho apresentou o desenvolvimento de uma plataforma WebAR de teste baseada em AR.js, utilizada para validar de forma padronizada os modelos gerados pelas diferentes ferramentas. Essa plataforma permitiu avaliar a viabilidade real de uso dos modelos em cenários educacionais, indo além de análises puramente teóricas ou visuais. Além disso, o protótipo desenvolvido evidencia a viabilidade técnica de integrar etapas fundamentais do fluxo de autoria — geração, validação e publicação — em uma solução baseada na Web.

Por fim, os resultados indicam que a principal lacuna atual não está na ausência de ferramentas de geração 3D, mas na falta de uma plataforma integrada voltada a usuários não técnicos, capaz de unificar a criação de avatares, a geração de marcadores, a inserção de áudio e interações, e a publicação automática em WebAR. Como trabalhos futuros, pretende-se expandir a plataforma desenvolvida neste estudo, incorporando funcionalidades de geração automática de avatares, criação de marcadores personalizados e configuração de fala e interações de forma no-code. Essa evolução pode contribuir significativamente para a democratização da autoria de experiências imersivas educacionais, ampliando o acesso à Realidade Aumentada como recurso pedagógico.

\bibliographystyle{IEEEtran}
\bibliography{referencia}

\end{document}
